% ID: SEL-D0007-BC-E05
% ESTADO: cotejada. Segunda pasada visual completa y borrador compilado revisados.
% Fuente: Selectividad/matematicas_aplicadas/ORD_TITULAR_MATEMÁTICAS APLICADAS CCSS-JUNIO-IMPRESO.pdf
% SHA-256: c42a605501a277fe7194b964ad92ad66047be9a62dc254fa2219d756d3d71317
% Página: 2; bloque C; ejercicio original 5.

Se desea probar la eficacia de dos tipos de vacunas, $A$ y $B$, contra un
virus determinado. Para ello, se seleccionan 5000 voluntarios sin anticuerpos
para este virus, a los que se les administra una de las vacunas o un placebo,
resultando que 3000 reciben la vacuna $A$, 1500 la $B$ y el resto el placebo.
Se comprueba que el 90\,\% de los vacunados con la $A$ y el 95\,\% de los
vacunados con la $B$, generan anticuerpos, no generando anticuerpos los que
han recibido el placebo. Se selecciona uno de esos voluntarios al azar.
\begin{enumerate}[label=\alph*)]
  \item \textbf{(1.5 puntos)} \textquestiondown{}Cuál es la probabilidad de
  que haya generado anticuerpos?
  \item \textbf{(1 punto)} Si dicho voluntario no ha generado anticuerpos,
  \textquestiondown{}qué probabilidad hay de que se le haya administrado
  placebo?
\end{enumerate}
